\chapter{DIFFERENTIALS (ADVANCED AND DESIRABLE SKILLS)}

This block focuses on advanced competencies that enhance performance in machine learning, generative AI, and large-scale experimentation environments.

\begin{center}
\textbf{These skills are not strictly required, but they significantly increase efficiency, scalability, and impact in real-world AI systems}
\end{center}

\vspace{0.5cm}

\section*{1. ADVANCED GENERATIVE AI}

\subsection*{1.1 Reinforcement Learning from Human Feedback (RLHF)}

RLHF is a technique used to align machine learning models with human preferences.

\subsubsection*{Pipeline Structure}

\begin{enumerate}
\item Pretrain a base model
\item Collect human feedback
\item Train a reward model
\item Optimize the policy using reinforcement learning
\end{enumerate}

\subsubsection*{Mathematical Formulation}

The objective is to maximize expected reward:

\[
\max_{\pi_\theta} \mathbb{E}_{x \sim \pi_\theta} [R(x)]
\]

where:
\begin{itemize}
\item $\pi_\theta$ is the model policy
\item $R(x)$ is the learned reward function
\end{itemize}

Policy update (example using policy gradient):

\[
\nabla_\theta J(\theta) = \mathbb{E} \left[ \nabla_\theta \log \pi_\theta(x) \cdot R(x) \right]
\]

\subsection*{1.2 Prompt Engineering and Dataset Design}

\subsubsection*{Prompt Engineering}

\begin{itemize}
\item Designing structured inputs to guide model outputs
\item Controlling behavior via instructions and context
\item Few-shot and zero-shot prompting
\end{itemize}

\subsubsection*{Dataset Engineering}

\begin{itemize}
\item Data curation and cleaning
\item Labeling and annotation strategies
\item Balancing and augmentation
\end{itemize}

Example of structured prompt:

\begin{minted}{python}
prompt = """
Task: Classify the sentiment of the sentence.
Input: The movie was excellent.
Output:
"""
\end{minted}

\vspace{0.5cm}

\section*{2. INFRASTRUCTURE}

\subsection*{2.1 GPU Computing}

GPU computing accelerates large-scale numerical operations:

\begin{itemize}
\item Parallel processing of tensors
\item Efficient training of deep neural networks
\item Support via CUDA and specialized libraries
\end{itemize}

Example:

\begin{minted}{python}
import torch

device = torch.device("cuda" if torch.cuda.is_available() else "cpu")
x = torch.randn(1000, 1000).to(device)
y = torch.matmul(x, x)
\end{minted}

\subsection*{2.2 Experimental Pipelines}

Structured workflows for running experiments:

\begin{itemize}
\item Data preprocessing
\item Model training
\item Evaluation and logging
\item Reproducibility
\end{itemize}

Pipeline abstraction:

\[
\text{Data} \rightarrow \text{Model} \rightarrow \text{Training} \rightarrow \text{Evaluation}
\]

\vspace{0.5cm}

\section*{3. MLOPS}

\subsection*{3.1 Overview}

MLOps (Machine Learning Operations) focuses on managing the lifecycle of machine learning models.

\subsection*{3.2 MLflow}

\begin{itemize}
\item Experiment tracking
\item Model versioning
\item Deployment support
\end{itemize}

Example:

\begin{minted}{python}
import mlflow

with mlflow.start_run():
    mlflow.log_param("lr", 0.01)
    mlflow.log_metric("loss", 0.25)
\end{minted}

\subsection*{3.3 Weights \& Biases}

\begin{itemize}
\item Real-time experiment tracking
\item Visualization dashboards
\item Hyperparameter tuning
\end{itemize}

Example:

\begin{minted}{python}
import wandb

wandb.init(project="example")
wandb.log({"loss": 0.25})
\end{minted}

\subsection*{3.4 DVC (Data Version Control)}

\begin{itemize}
\item Versioning datasets and models
\item Reproducible pipelines
\item Integration with Git
\end{itemize}

Example:

\begin{minted}{python}
# Command-line usage (conceptual)
# dvc init
# dvc add data.csv
# git add data.csv.dvc
\end{minted}

\vspace{0.5cm}

\section*{FINAL CONNECTION}

These advanced skills enhance:

\begin{enumerate}
\item \textbf{Model Alignment:} RLHF and prompt engineering
\item \textbf{Scalability:} GPU computing and pipelines
\item \textbf{Reproducibility:} MLOps tools
\item \textbf{Deployment Readiness:} experiment tracking and versioning
\end{enumerate}

\vspace{0.5cm}

\section*{STUDY ROADMAP}

\subsection*{Phase 1 --- Generative AI}
\begin{itemize}
\item RLHF fundamentals
\item Prompt engineering techniques
\end{itemize}

\subsection*{Phase 2 --- Infrastructure}
\begin{itemize}
\item GPU acceleration
\item Experiment pipelines
\end{itemize}

\subsection*{Phase 3 --- MLOps}
\begin{itemize}
\item MLflow
\item Weights \& Biases
\item DVC
\end{itemize}

\subsection*{Phase 4 --- Integration}
\begin{itemize}
\item End-to-end ML workflows
\item Scalable and reproducible systems
\end{itemize}